\documentclass[a4paper,11pt]{report}
\usepackage[utf8]{inputenc}
\usepackage[francais]{babel}
\usepackage{graphicx}


% Title Page
\title{SPAM ASSASSIN}
\date{Décembre 2014}
\author{Mathieu KERN -}


\begin{document}
\maketitle
\includegraphics{spamassassinintro.png}
\tableofcontents

\pagebreak

\part{Présentation de SpamAssassin}

\section{Problématique }

  Le mail(ou courriel) est aujourd'hui le moyen privilégié de communication à travers le monde. Massivement utilisé, 
d'une certaine fiabilité et éprouvé par des décennies d'utilisation il reste le moyen le plus répandues pour 
les communications entres les personnes. Malheureusement, mail est également aujourd'hui synonyme de spam, ces messages
indésirables qui s'entassent dans nos boites mails. C'est ici qu'entre en jeu SpamAssassin.

\subsection{Le SPAM}
Avant de poursuivre sur SpamAssassin, rappellons concrètement ce qu'est le SPAM et ce qu'il implique. 

\subsubsection{Ses caractéristiques}

\paragraph{Comment reconnaître un SPAM:}

\begin{itemize}
 \item De par sa nature, un SPAM n'est pas désiré par l'utilisateur qui le reçoit. 
 \item La réception d'un SPAM résulte d'un envoi massif: une machine (souvent un bot) envoi le même message 
 à plusieurs destinataire sans aucun discernement. Cela s'oppose aux messages ciblés par exemple de commerçant,
 qui n’envoie que à leurs prospects.
 \item Son contenu n'est pas destiné spécifique à l'utilisateur (chaque personne reçoit le même contenu)
 \item Une importante liste de destinataires 
 \item Entête des messages souvent corrompues ou ne respectant pas les normes
\end{itemize}

\paragraph{Statut légal}


La loi pour la confiance dans l'économie numérique du 21 juin 2004 contient une transposition de la
directive européenne du 12 juillet 2002\footnote{Le principe introduit figura également à l'article L.34.5 du code 
des postes et des communications électroniques } relative à la protection de la vie privée dans le secteur des communications
électroniques:
\begin{quote}
 Est interdite la prospection directe au moyen d'un automate d'appel, d'un télécopieur ou d'un courrier électronique utilisant,
 sous quelque forme que ce soit, les coordonnées d'une personne physique qui n'a pas exprimé son consentement préalable à recevoir
 des prospections directes par ce moyen. 
\end{quote}
Les SPAM sont donc connus du droit français et encadrés par des textes spécifiques.

\

\subsection{Le principe de Spam Assassin}

\pagebreak

\appendix





\end{document}          
